\documentclass{article}
\usepackage[T1]{fontenc}
\usepackage{lmodern}
\usepackage{graphicx}
\usepackage{amsmath,amssymb}
\usepackage[a4paper,margin=2cm]{geometry}
\usepackage[absolute,overlay]{textpos}
\setlength{\TPHorizModule}{1mm}
\setlength{\TPVertModule}{1mm}
\usepackage[indonesian]{babel}
\usepackage{hyperref}
\setlength{\emergencystretch}{2em}
% unit: Halaman 16

\begin{document}

% === Halaman 16 (python/native) ===

16

• Karena setiap blok plainteks yang sama selalu dienkripsi menjadi blok cipherteks yang 

sama, maka secara teoritis dimungkinkan membuat buku kode plainteks dan cipherteks 

yang berkoresponden (asal kata “code book” di dalam ECB ). Enkripsi/dekripsi dilakukan 

dengan look up table. 

        Plainteks            Cipherteks


0000 


0100


0001 


1001 


0010 


1010


… 


…


1111 


1010 


• Untuk setiap kunci K yang berbeda, dibuat buku kode yang berbeda pula. Jika panjang 

kunci n bit, maka terdapat 2n buku kode.

\end{document}
